% ============================================================================
%  Task_W01.tex -- Recursion Basics, Week 1
%  Uses coa-accessible.sty for the accessible baseline (headings, captions,
%  color palette, page numbers, PDF metadata). Compile twice with pdflatex
%  so "Page X of Y" resolves.
% ============================================================================
\documentclass[11pt]{article}
\usepackage{amsmath}
\usepackage{coa-accessible}

\AssignmentMeta{Task W01: Recursion Basics}{Week 1 (Base Case)}

\begin{document}

\AssignmentTitleBlock

\begin{center}
\rule{0.95\linewidth}{0.6pt}
\end{center}
\subsection*{A note on how this assignment was developed}
\noindent
This assignment was drafted with AI help (mostly for brainstorming new
problems to solve and formatting with LaTeX). It should be accurate, but
something might still be wrong.
Because this is a new assignment, you are also helping me debug it, and I
will take that into consideration when I review your work. That said,
helping me find rough edges is not a substitute for doing the assignment;
each problem below is solvable, but you may need to make assumptions in
your attempt. If you find an error or
something unclear, please tell me; that helps later versions, though
finding errors is not the point of this assignment.
\begin{center}
\rule{0.95\linewidth}{0.6pt}
\end{center}
\vspace{0.6em}

\section{Important Guidelines}
\begin{itemize}[leftmargin=1.4em]
  \item Exam questions and problems may extend beyond tasks and lab assignments.
  \item No Generative AI: You must complete this assignment using only the
        resources provided in the course (book, lectures, and materials).
        Credit will not be granted for using external features or techniques
        not covered in this course. (See the one exception in the Applied
        Problems section.)
  \item Ask Me Questions: If you need clarification, please ask me directly.
        Other instructors or advisors may not be aware of the specific
        details and expectations of this assignment.
  \item The assignment instructions do not detail every step; therefore, you
        must read, interpret, and possibly make assumptions. If something is
        unclear, write down the assumption you make, proceed with it, and
        complete the problem. Do not leave a problem unfinished because a
        detail was ambiguous. Not all assumptions are automatically correct,
        so state yours clearly enough that I can follow your reasoning.
\end{itemize}

\section{Diagram Submission Instructions}
You are required to create and submit diagrams for this assignment. You may
choose one of the following options.

\begin{itemize}[leftmargin=1.4em]
  \item \textbf{Hand-drawn on paper or a whiteboard.} Scan or take clear,
        well-lit photos. Images must be easy to read and clearly labeled
        with problem numbers. This is often quicker than using diagram
        tools.
  \item \textbf{Hand-drawn using an app} (for example Vittle Pro or Adobe
        Fresco). Save as an image file (PNG, JPG) or PDF. This is a
        convenient option but may be harder without a stylus.
  \item \textbf{Created with diagramming software} (for example Draw.io,
        LibreOffice Draw, or Excalidraw). Save in the tool's native file
        format; do not export as an image or PDF.
  \item \textbf{Other tools.} If you would like to use a different
        diagramming application, ask for approval first. Note that I have
        limited installation access on my office computer.
\end{itemize}

\noindent Place one problem per tab. Include your name and the problem
number on every diagram.

\subsection*{If you find something that needs to be debugged}
If, while working through this assignment, you discover something in the
problem statements that needs to be debugged, describe what you found in
the "Bugs or Issues You Found" section of the answer sheet.

\subsection*{Where to submit}
Commit your full submission, including your diagrams, your pseudocode,
your answer sheet, and your AI conversation if you used the formatting
exception, to a folder named \texttt{taskw01} in your git repository.

\section{Recurrence Relation Problems}

Both problems below are required. For each one, draw the full recursion
tree for the given value of $n$, and record the final value the relation
returns on your answer sheet.

\subsection{Problem: $B(n)$}
\[
B(n) = \begin{cases}
1, & n \le 2 \\
B(n-1) - B(n-2), & n > 2
\end{cases}
\]
Draw the recursion tree for $B(6)$. Record the value of $B(6)$ on your
answer sheet.

\subsection{Problem: $D(n)$}
\[
D(n) = \begin{cases}
1, & n \le 3 \\
D(n-2) + D(n-1) - D(n-3), & n > 3
\end{cases}
\]
Draw the recursion tree for $D(6)$. Record the value of $D(6)$ on your
answer sheet.

\section{Applied Problems}

Four problems are described below, in no particular order of difficulty.
Read through all four before you start.

If you have not written pseudocode before, or would like to see it written
the way Cormen, Leiserson, Rivest, and Stein (CLRS) and Carrano's Walls and
Mirrors write it, this guide covers those conventions with worked examples:
\url{https://course.ccs.neu.edu/cs3000/resources/latex_pseudocode.pdf}. You
do not need to use the LaTeX package described in that guide, or match its
typesetting exactly; plain pseudocode that follows the same conventions
(numbered lines, indentation, real \texttt{for}/\texttt{while}/\texttt{if}
structure) is fine. You may also write your solution as working code in
Java or C/C++ instead of pseudocode, if you would rather work that way.

Once you have solved a problem, and only after you have solved it, you may
use an AI tool solely to format your finished pseudocode into the
conventions above. If you do this, you must explicitly tell the AI not to
change or fix your algorithm, only to format it, and you must submit the
full AI conversation alongside your pseudocode. This exception does not
extend to using AI to solve the problem, or to fix a mistake in your
algorithm; that is still covered by the No Generative AI guideline above.

\begin{itemize}[leftmargin=1.4em]
  \item Every solution must use recursion. A solution that does not use
        recursion will not receive credit, even if it produces the correct
        result.
  \item Choose \textbf{one} of the four and complete it in full. Write
        correct recursive pseudocode, or working recursive code in Java or
        C/C++, for a solution, and draw a full diagram tracing your
        solution on the example input given. A box trace or a recursion
        tree is acceptable, whichever fits your solution. Record the final
        returned value on your answer sheet.
  \item For the \textbf{other three}, write your best attempt at a
        recursive solution, in pseudocode or in Java or C/C++. No diagram
        is required for these three.
\end{itemize}

\subsection{Problem 1}
Write a method that takes a positive integer and returns the sum of its
digits. For example, the digit sum of 48173 is $4+8+1+7+3 = 23$.

\vspace{4pt}\noindent\textit{Example input if you choose this one:} 48173.

\subsection{Problem 2}
A staircase has $n$ steps. You can climb it by taking either 1 step or 2
steps at a time. Write a method that returns the number of different ways
to climb a staircase with $n$ steps. A recursion tree is an acceptable
diagram for this problem, if a box trace does not fit your solution.

\vspace{4pt}\noindent\textit{Example input if you choose this one:} 6 steps.

\subsection{Problem 3}
Write a method that takes a string of characters and returns true if the
string reads the same forwards and backwards, and false otherwise.

\vspace{4pt}\noindent\textit{Example input if you choose this one:} \texttt{"racecar"}.

\subsection{Problem 4}
Write a method that returns the greatest common divisor of two positive
integers $a$ and $b$, using Euclid's algorithm: the greatest common divisor
of $a$ and $b$ is the same as the greatest common divisor of $b$ and the
remainder of $a$ divided by $b$, and the greatest common divisor of any
number and 0 is that number.

\vspace{4pt}\noindent\textit{Example input if you choose this one:} $a = 150$, $b = 84$.

\end{document}
